% ===== PRÉAMBULE CANONIQUE — à coller VERBATIM en tête de CHAQUE .tex =====
\documentclass[11pt,a4paper]{article}
\usepackage[utf8]{inputenc}\usepackage[T1]{fontenc}\usepackage{lmodern}
\usepackage[french]{babel}
\usepackage{amsmath,amssymb,mathtools}\usepackage{icomma}
\usepackage{siunitx}\sisetup{locale=FR}  % \qty{}{}, \si{}, \ang{}
\usepackage{xcolor}\usepackage{colortbl}
\usepackage{tikz}\usepackage{pgfplots}\pgfplotsset{compat=1.18}
\usepackage[siunitx,europeanresistors]{circuitikz} % circuits (élec.)
\usepackage[most]{tcolorbox}\usepackage{enumitem}\usepackage{array,booktabs}
\usepackage[a4paper,margin=2.2cm]{geometry}
\usetikzlibrary{arrows.meta,calc,angles,quotes,positioning,patterns,%
  backgrounds,shapes.geometric,decorations.pathmorphing,decorations.markings,intersections}
\definecolor{primary}{RGB}{222,49,99}\definecolor{primarydark}{RGB}{150,22,55}
\definecolor{defcol}{RGB}{0,114,178}\definecolor{methcol}{RGB}{52,92,148}
\definecolor{excol}{RGB}{0,135,90}\definecolor{attncol}{RGB}{200,40,55}
\tcbset{boxbase/.style={enhanced,breakable,boxrule=0.9pt,arc=2.5pt,
  left=3mm,right=3mm,top=2.3mm,bottom=2.3mm,fonttitle=\bfseries,coltitle=white}}
% --- boîtes TUTORIEL (noms EXACTS, ne pas renommer) ---
\newtcolorbox{cle}[1][]{boxbase,colback=defcol!6,colframe=defcol,
  title={\textsc{À retenir}\ifx\relax#1\relax\relax\else\textmd{ : \itshape #1}\fi}}
\newtcolorbox{methode}[1][]{boxbase,colback=methcol!7,colframe=methcol,sharp corners,
  title={\textsc{Méthode}\ifx\relax#1\relax\relax\else\textmd{ --- \itshape #1}\fi}}
\newtcolorbox{exemplebox}[1][]{boxbase,colback=excol!5,colframe=excol,
  title={\textsc{Exemple}\ifx\relax#1\relax\relax\else\textmd{ : \itshape #1}\fi}}
\newtcolorbox{piege}[1][]{boxbase,colback=attncol!6,colframe=attncol,
  title={\textsc{Piège}\ifx\relax#1\relax\relax\else\textmd{ : \itshape #1}\fi}}
\newtcolorbox{remarque}[1][]{boxbase,colback=black!4,colframe=black!45,
  title={\textsc{Remarque}\ifx\relax#1\relax\relax\else\textmd{ : \itshape #1}\fi}}
% --- boîtes EXERCICE / CORRIGÉ (noms EXACTS, auto-comptées) ---
\newtcolorbox[auto counter]{exo}[1][]{boxbase,colback=primary!4,colframe=primary,
  title={\textsc{Exercice}~\thetcbcounter\ifx\relax#1\relax\relax\else\textmd{ --- \itshape #1}\fi}}
\newtcolorbox[auto counter]{corrige}[1][]{boxbase,colback=excol!4,colframe=excol,
  title={\textsc{Corrigé}~\thetcbcounter\ifx\relax#1\relax\relax\else\textmd{ --- \itshape #1}\fi}}
\newcommand{\vv}[1]{\vec{#1}}\newcommand{\ds}{\displaystyle}
\newcommand{\R}{\mathbb{R}}\newcommand{\N}{\mathbb{N}}\newcommand{\Z}{\mathbb{Z}}
% (un \usepackage{fancyhdr}+en-tête est possible ; il est ignoré côté web)
% =========================================================================
\begin{document}
\sisetup{per-mode=symbol}

\begin{corrige}[hauteur maximale avec vitesse initiale]
\begin{enumerate}[label=\alph*)]
  \item Le sac part avec la vitesse ascensionnelle de la montgolfière ($v_0=\qty{10}{\metre\per\second}$ vers le haut). Il monte encore de $\dfrac{v_0^2}{2g}=\dfrac{100}{20}=\qty{5}{\metre}$ : hauteur maximale $=120+5=\qty{125}{\metre}$.
  \item Même raisonnement depuis \qty{3}{\metre} : $3+\dfrac{10^2}{2\times10}=3+5=\qty{8}{\metre}$.
\end{enumerate}
\end{corrige}

\begin{corrige}[quand la trajectoire fait 45°]
\begin{enumerate}[label=\alph*)]
  \item À \ang{45}, $\tan\ang{45}=1$ donc $v_y=v_x$. Or $v_x=v_{0x}=\qty{40}{\metre\per\second}$ (constant) et $v_y=g\,t$ : $10\,t=40 \Rightarrow t=\qty{4}{\second}$.
  \item Chute : $y=\tfrac12 g t^2=\tfrac12\times10\times16=\qty{80}{\metre}$. Vitesse : $v=\sqrt{v_x^2+v_y^2}=\sqrt{40^2+40^2}=40\sqrt2\approx\qty{56,6}{\metre\per\second}$.
\end{enumerate}
\end{corrige}

\begin{corrige}[la masse n'y change rien]
\begin{enumerate}[label=\alph*)]
  \item \textbf{Oui, même hauteur au même instant.} En chute libre, l'accélération $g$ est indépendante de la masse : $y_{\max}=v_0^2/2g$ et le temps de vol $2v_0/g$ ne dépendent pas de $m$.
  \item $y_{\max}=\dfrac{v_0^2}{2g}=\dfrac{20^2}{2\times1{,}6}=\dfrac{400}{3{,}2}=\qty{125}{\metre}$ (bien plus haut que sur Terre, car $g$ y est six fois plus faible).
\end{enumerate}
\end{corrige}

\begin{corrige}[la profondeur du puits]
\begin{enumerate}[label=\alph*)]
  \item $t_c+t_s=\qty{4,25}{\second}$, avec $h=\tfrac12 g\,t_c^2$ (chute) et $t_s=\dfrac{h}{v_{\text{son}}}$ (remontée du son).
  \item Si $t_c=\qty{4}{\second}$ : $h=\tfrac12\times10\times16=\qty{80}{\metre}$, d'où $t_s=\dfrac{80}{320}=\qty{0,25}{\second}$. Total $4+0{,}25=\qty{4,25}{\second}$ $\checkmark$. La profondeur est de \qty{80}{\metre}.
\end{enumerate}
\end{corrige}

\begin{corrige}[lancer vertical : vrai ou faux]
\begin{enumerate}[label=\alph*)]
  \item \textbf{V} : $y_{\max}=\dfrac{v_0^2}{2g}=\dfrac{100}{20}=\qty{5}{\metre}$.
  \item \textbf{F} : le temps de montée est $t_s=v_0/g=\qty{1}{\second}$, mais la balle ne retombe au sol qu'après $2t_s=\qty{2}{\second}$.
  \item \textbf{V} : au sommet la vitesse verticale s'annule ($v=v_0-gt=0$).
  \item \textbf{V} : par symétrie de la chute libre, elle repasse au sol à \qty{10}{\metre\per\second} (sens opposé).
\end{enumerate}
\textbf{Bilan : 3 affirmations vraies} (a, c, d).
\end{corrige}

\end{document}
